\section{Conclusion}

On my laptop (6 cores, 12 logical processors) all the work I measured, counting 50 million
numbers, was \textbf{CPU-bound}. I tell this by two checks: the time falls when threads are added
(as long as there are free cores), and the time is close to the time of the pure computation, with
no part spent waiting for a device.
\begin{itemize}
\item \emph{Task 1a}: the time fell from 919\,ms on 1 thread to 301\,ms at 6 threads and 279\,ms at
12. If the work were waiting for the disk, more threads would not have made it three times faster.
Beyond the 12 logical processors it got worse, because there was no more processor to give.
\item \emph{Task 1b}: the data is already in memory, so there is no I/O at all. The per-thread
version went from 182\,ms on 1 thread to 68\,ms at 8 threads.
\item \emph{Task 1c} is limited by the processor too, but by only \emph{one} core. The single
reader that reads the lines and creates the strings is the bottleneck: the workers find the queue
empty for almost every batch, and 1 to 64 workers all take about 2.1\,s.
\end{itemize}
None of these tasks was I/O-bound on my machine. The 369\,MB file fits easily in the 16\,GB of RAM,
so Windows keeps it in its file cache, and ``reading'' it is only copying memory. Work becomes
I/O-bound when the data has to come from a device much slower than the processor, which these
tasks did not do.

Synchronisation has a price that depends on how often it happens. A lock taken for every number
made Task 1b ten times slower than per-thread counters (about 770\,ms vs.\ 75\,ms) and never faster
than one thread. Batches of one line made Task 1c 6--20 times slower than batches of 10\,000. Shared
counters without synchronisation were both wrong (84--87\,\% of the increments lost) and over
4 times slower than per-thread counters (about 330\,ms). The fastest correct design was to share
nothing while working and combine the results once at the end.
